% !PRINTED 587
\weberpage{587}

Ist $\alpha$ durch $\beta$ theilbar, so ist $N(\alpha)$ durch $N(\beta)$ theilbar.
Denn aus $\alpha=\beta\gamma$ folgt $N(\alpha)=N(\beta)N(\gamma)$.

Sind $\alpha$, $\beta$, $\gamma$ ganze Zahlen in $R(i)$ und ist $\alpha$ durch $\gamma$ theilbar,
so ist auch $\alpha\beta$ durch $\gamma\beta$ theilbar, und ist $\alpha$ und $\beta$ durch $\gamma$
theilbar, so ist auch $\alpha\pm\beta$ durch $\gamma$ theilbar.

Jede ganze Zahl ist durch jede Einheit theilbar.

% [source-original mathematical ellipsis, verified 2026-08-29] From
% $N(\alpha)N(\beta)=1$ with positive integral norms, both norms equal $1$;
% the print omits ``$=1$'' after the first $N(\alpha)$.
Die Einheiten sind aber nur durch Einheiten theilbar; denn
ist $\alpha\beta$ eine Einheit, so ist $N(\alpha)N(\beta)=1$, also $N(\alpha)$ und
$N(\beta)=1$, d.\,h. $\alpha$ und $\beta$ sind Einheiten.

Associirte Zahlen sind gegenseitig durch einander theilbar,
und wenn zwei Zahlen gegenseitig theilbar sind, so sind sie
associirt. Denn sind $\alpha:\beta$ und $\beta:\alpha$ beides ganze Zahlen, so ist
das Product dieser beiden Zahlen $=1$; also sind beides Einheiten.

Für die Zahlen des Körpers $R(i)$ gelten dieselben
Gesetze über die Zerlegung in Primfactoren, wie bei den reellen
ganzen Zahlen.

\gesperrt{Eine ganze Zahl $\alpha$ des Körpers $R(i)$, die keine Einheit
 ist, heisst zusammengesetzt, wenn sie sich in mehrere
ganzzahlige Factoren, deren keiner eine Einheit ist, zerlegen
lässt.}

\gesperrt{Lässt sie sich nicht so zerlegen, so soll sie eine
Primzahl im Körper $R(i)$ heissen.}

Ist $\xi=x+yi$ eine gebrochene Zahl des Körpers $R(i)$, so
lässt sich eine ganze Zahl $\mu=m+ni$ so bestimmen, dass die
Norm der Differenz $\xi-\mu$, also $(x-m)^2+(y-n)^2$, kleiner
oder wenigstens nicht grösser als $1/2$ ist; denn man braucht die
ganzen rationalen Zahlen $m$, $n$ nur so zu wählen, dass $x-m$
und $y-n$ absolut genommen nicht grösser als $1/2$, ihre Quadrate
also nicht grösser als $1/4$ sind.

Sind also $\alpha$ und $\alpha_1$ zwei ganze Zahlen in $R(i)$, und $\alpha_1$ von
Null verschieden, so kann man hiernach die ganzen Zahlen $\mu$
und $\alpha_2$, so bestimmen, dass
\[
\frac{\alpha}{\alpha_1}-\mu=\frac{\alpha_2}{\alpha_1},\quad \frac{N(\alpha_2)}{N(\alpha_1)}\leqq \frac{1}{2},
\]
also
\[
\alpha=\mu\alpha_1+\alpha_2,\quad N(\alpha_2)<N(\alpha_1).
\]
Ist $\alpha_2$ nicht Null, so kann man ebenso mit den Zahlen
$\alpha_1$, $\alpha_2$ verfahren und erhält
\[
\alpha_1=\mu_1\alpha_2+\alpha_3,\quad N(\alpha_3)<N(\alpha_2),
\]
%CONT>>
